\documentclass{article}

\usepackage{amsmath, amsfonts, amssymb, amsthm}
\usepackage{parskip}

\newtheorem*{claim}{Claim}

\title{Proof by Induction: Derivatives of $x e^x$}
\author{Introduction to Engineering Mathematics}
\date{November 2026}

\begin{document}
\maketitle

\begin{claim}
  Let $f(x) = x e^x$. For every integer $n \ge 1$,
  \[
    f^{(n)}(x) = n e^x + x e^x = (n + x) e^x.
  \]
\end{claim}

\section*{Proof outline}

Let $P(n)$ denote the statement ``$f^{(n)}(x) = n e^x + x e^x$''.
We show that $P(n)$ holds for all $n \ge 1$ by induction on $n$.

\subsection*{Step 1: Base case ($n = 1$)}

Compute the first derivative using the product rule:
\[
  f'(x) = \frac{d}{dx}(x) \cdot e^x + x \cdot \frac{d}{dx}(e^x)
        = e^x + x e^x.
\]
This is $1 \cdot e^x + x e^x$, so $P(1)$ holds.

\subsection*{Step 2: Induction hypothesis}

Assume that $P(k)$ holds for some integer $k \ge 1$, that is,
\[
  f^{(k)}(x) = k e^x + x e^x.
\]

\subsection*{Step 3: Induction step}

We must show that $P(k + 1)$ holds, i.e.\ that
$f^{(k+1)}(x) = (k + 1) e^x + x e^x$.

The $(k+1)$-th derivative is the derivative of the $k$-th derivative, so by
the induction hypothesis
\[
  f^{(k+1)}(x) = \frac{d}{dx} f^{(k)}(x)
               = \frac{d}{dx} \left( k e^x + x e^x \right).
\]
Differentiate each term:
\begin{itemize}
\item $\dfrac{d}{dx}(k e^x) = k e^x$, since $k$ is a constant;
\item $\dfrac{d}{dx}(x e^x) = e^x + x e^x$, by the product rule (as in the
  base case).
\end{itemize}
Adding these together gives
\[
  f^{(k+1)}(x) = k e^x + e^x + x e^x = (k + 1) e^x + x e^x,
\]
which is exactly $P(k + 1)$.

\subsection*{Step 4: Conclusion}

Since $P(1)$ is true, and $P(k)$ implies $P(k + 1)$ for every $k \ge 1$, the
principle of mathematical induction tells us that $P(n)$ is true for all
$n \ge 1$. \qed

\end{document}
